\documentclass[12pt,a4paper]{article}
% Поддержка русского языка
\usepackage{polyglossia}
\setmainlanguage{russian}
\setotherlanguage{english}
\setmainfont{Times New Roman}
\newfontfamily{\cyrillicfont}{Times New Roman}
\newfontfamily{\cyrillicfontsf}{Arial}
\newfontfamily{\cyrillicfonttt}{Courier New}

\usepackage{amsmath,amssymb,amsthm,mathtools}
\usepackage{geometry}
\geometry{left=1.5cm,right=1.5cm,top=1.5cm,bottom=2.0cm}
\usepackage{booktabs}
\usepackage{longtable}
\usepackage{array}
\usepackage{hyperref}
\usepackage{url}
\usepackage{xcolor}
\usepackage{titlesec}
\usepackage{fancyhdr}
\usepackage{enumitem}
\usepackage{makecell}
\usepackage{float}

\titleformat{\section}{\Large\bfseries}{\thesection}{1em}{}
\titleformat{\subsection}{\large\bfseries}{\thesubsection}{1em}{}

\hypersetup{colorlinks=true,linkcolor=blue,citecolor=blue,urlcolor=blue}

% --- Фундаментальные константы YUCT (оставляем как есть) ---
\newcommand{\REL}{\mathcal{K}}          % относительная эффективность координации
\newcommand{\ABS}{K_{\mathrm{eff}}}     % абсолютная
\newcommand{\kappac}{\kappa_c}
\newcommand{\betaf}{\beta}
\newcommand{\Sodd}{S_{\mathrm{odd}}}
\newcommand{\Seven}{S_{\mathrm{even}}}

\begin{document}
	
	\begin{titlepage}
		\centering
		\vspace*{2cm}
		{\Huge\bfseries Приложение BCLM-LL\\[0.3cm]
			Экспериментальный протокол для проверки\\долгоживущего координационного дизайна в культивируемых кардиомиоцитах человека\par}
		\vspace{0.5cm}
		{\Large\bfseries Фальсифицируемая проверка координационной теории старения YUCT\par}
		\vspace{1.5cm}
		{\large A.\,V.\,Якушев\par}
		{\normalsize Исследования YUCT, \url{https://yuct.org/}\par}
		\vspace{0.5cm}
		{\normalsize Главный DOI: \url{https://doi.org/10.5281/zenodo.18444598}\par}
		\vspace{0.5cm}
		{\normalsize Июнь 2026 (v1.2 – исправлена численная согласованность)\par}
		\vfill
		\begin{abstract}
			\noindent
			Объединённая теория координации Якушева (YUCT) объясняет старение как прогрессирующее снижение эффективности координации \(K_{\mathrm{eff}}\) клеточных подсистем. Биологический лагранжиан координации (BCLM, Приложение BCLM) предоставляет математическую основу для вычисления \(K_{\mathrm{eff}}\) по последовательностям и для конструирования генетических вмешательств, восстанавливающих её до ювенильного уровня. Данный протокол описывает полный, самодостаточный эксперимент по проверке долгоживущего координационного дизайна в культивируемых кардиомиоцитах человека. Мы указываем точные гены, подлежащие модификации, вычисляем их нативную и оптимизированную относительную эффективность координации \(\REL\), выводим количественные предсказания для частоты мутаций, окислительного стресса, агрегации белков, скорости миграции и клеточной продолжительности жизни непосредственно из универсального закона ошибок \(\varepsilon = \kappa_c \alpha \REL^{-\beta}\) (\(\beta=2/3\), \(\kappa_c=1/3\)), и предоставляем пошаговый экспериментальный график. Включён модуль контролируемой регрессии (CCR) для демонстрации обратимости стареющего фенотипа. Все предсказания предварительно зарегистрированы и фальсифицируемы с помощью стандартных методов клеточной биологии. По этическим соображениям протокол намеренно опускает методы \textit{in vivo} доставки; читатель отсылается к Приложению~\(\Sigma\).
		\end{abstract}
	\end{titlepage}
	
	\tableofcontents
	\newpage
	
	\section{Введение}
	\label{sec:intro}
	Универсальный закон масштабирования ошибок YUCT,
	\begin{equation}
		\varepsilon = \kappa_c \, \alpha \, \REL^{-\beta}, \qquad \beta = \frac{2}{3},\quad \kappa_c = \frac{1}{3},
		\label{eq:error-law}
	\end{equation}
	утверждает, что относительная ошибка \(\varepsilon\) любого координированного процесса (трансляция, репарация, перенос электронов) зависит только от относительной эффективности координации \(\REL = K_{\mathrm{eff}}/K_{\mathrm{eff}}^{\max}\) и системно-специфической константы \(\alpha\) порядка \(10^{-2}\)–\(10^{0}\). Этот закон был подтверждён на более чем 40 порядках величины в физических, химических и биологических системах (см. Приложение L и сводную валидационную заметку).
	
	В молодой клетке все многобелковые машины работают с высоким \(\REL\) (обычно 0.5–0.9 в относительных единицах). С возрастом накопленные молекулярные повреждения снижают \(\REL\). Поскольку ошибки порождают дальнейшие повреждения, возникает порочный круг: \(\REL\) падает, ошибки растут, и клетка в конечном итоге входит в сенесценцию или апоптоз.
	
	Биологический лагранжиан координации (BCLM, Приложение BCLM) даёт единое операциональное определение \(\REL\) для кодонов, белков и ферментов, а также набор правил для предсказания того, как генетические изменения изменяют \(\REL\). Используя BCLM, можно \textbf{сконструировать} клетку, чей \(\REL\) зафиксирован на ювенильном уровне. Настоящий протокол превращает этот дизайн в конкретный, фальсифицируемый эксперимент.
	
	\subsection{Как BCLM вычисляет \(\REL\) для белка (краткое напоминание)}
	Для белка длины \(N\) абсолютная эффективность координации равна
	\[
	K_{\mathrm{eff}}^{\mathrm{protein}} = \Bigl( \prod_{i=1}^{N} K_{\mathrm{eff},i}^{\mathrm{AA}} \Bigr)^{1/N} \cdot R_{\mathrm{struct}},
	\]
	где \(K_{\mathrm{eff},i}^{\mathrm{AA}}\) — аминокислотно-специфические значения (Таблица 2 BCLM), а \(R_{\mathrm{struct}}\) — структурный резонансный фактор (Таблица 10 BCLM). Относительная эффективность \(\REL^{\mathrm{protein}} = K_{\mathrm{eff}}^{\mathrm{protein}}/K_{\mathrm{eff}}^{\max}\).
	
	Оптимизация кодонов повышает \(\REL\), поскольку каждый кодон имеет свою собственную \(\REL^{\mathrm{codon}}\) (Таблица 1 BCLM). Когда каждый кодон заменяется на его синонимичный аналог с наивысшим \(\REL\), среднее геометрическое эффективностей кодонов возрастает, и, следовательно, \(\REL\) всего белка увеличивается.
	
	\subsection{Температурная зависимость}
	Из Приложения P, \(K_{\mathrm{eff}}\) масштабируется с температурой как \(T^{-\gamma}\) (\(\gamma\approx0.3\)). Таким образом \(\varepsilon(T) \propto T^{0.20}\). Все эксперименты должны проводиться при строго контролируемой температуре 37.0\(^\circ\)C. Параллельная группа при 35\(^\circ\)C позволит проверить предсказание температурного масштабирования.
	
	\section{Долгоживущий дизайн: четыре целевых слоя}
	\label{sec:layers}
	Мы модифицируем четыре независимых координационных слоя. Таблица~\ref{tab:targets} перечисляет гены, их нативный и оптимизированный \(\REL\), а также индивидуальное изменение соответствующей функциональной ошибки. Эти числа вычислены с использованием конвейера BCLM (v6.1) с частотами кодонов человека из базы Kazusa и значениями \(R_{\mathrm{struct}}\) из Таблицы 10 BCLM.
	
	\begin{table}[H]
		\centering
		\caption{Целевые гены и предсказанные изменения \(\REL\).}
		\label{tab:targets}
		\begin{tabular}{@{}llrrr@{}}
			\toprule
			\textbf{Слой} & \textbf{Ген} & \(\REL_{\mathrm{WT}}\) & \(\REL_{\mathrm{LL}}\) & \(\varepsilon_{\mathrm{LL}}/\varepsilon_{\mathrm{WT}}\) \\
			\midrule
			1 (Геном) & BRCA1 & 0.62 & 0.78 & 0.85 \\
			& PARP1 & 0.60 & 0.77 & 0.86 \\
			& MLH1  & 0.59 & 0.76 & 0.87 \\
			2 (Митохондрии)   & NDUFS1 & 0.55 & 0.72 & 0.80 \\
			& NDUFV1 & 0.56 & 0.73 & 0.80 \\
			& NDUFA9 & 0.54 & 0.70 & 0.81 \\
			3 (Протеостаз) & HSPA1A & 0.52 & 0.69 & 0.78 \\
			& DNAJB1 & 0.50 & 0.68 & 0.79 \\
			& HSPA9  & 0.53 & 0.70 & 0.78 \\
			4 (ВКМ)    & ITGAV  & 0.58 & 0.74 & 0.83 \\
			& ITGB3  & 0.57 & 0.73 & 0.83 \\
			& FN1    & 0.60 & 0.76 & 0.84 \\
			\bottomrule
		\end{tabular}
		\footnotesize Отношение ошибок следует из \(\varepsilon_{\mathrm{LL}}/\varepsilon_{\mathrm{WT}} = (\REL_{\mathrm{LL}}/\REL_{\mathrm{WT}})^{-2/3}\).
	\end{table}
	
	\subsection{Слой 1 – Стабильность генома}
	Снижение частоты мутаций HPRT является комбинированным эффектом трёх оптимизированных репарационных белков. Поскольку пути репарации частично избыточны, общая вероятность ошибки приблизительно равна среднему геометрическому индивидуальных факторов:
	\[
	\frac{\mu_{\mathrm{LL}}}{\mu_{\mathrm{WT}}} \approx \sqrt[3]{0.85 \times 0.86 \times 0.87} \approx 0.86.
	\]
	Примечание: это ожидаемое снижение \emph{частоты точковых мутаций}. Произведение трёх факторов равно \(0.63\); среднее геометрическое учитывает перекрывающуюся репарационную активность и является более консервативным. Поэтому мы предсказываем, что частота мутаций снизится примерно до \(86\%\) от WT.
	
	\subsection{Слой 2 – Митохондриальный энергетический резонанс}
	Попарные совместимости \(C_{AB}\) были подняты выше 0.95. Предсказанное снижение продукции супероксида (флуоресценция MitoSOX) является средним геометрическим снижений ошибок отдельных субъединиц:
	\[
	\frac{\mathrm{ROS}_{\mathrm{LL}}}{\mathrm{ROS}_{\mathrm{WT}}} \approx \sqrt[3]{0.80 \times 0.80 \times 0.81} \approx 0.80.
	\]
	Таким образом, ожидается, что ROS упадёт примерно до \(80\%\) от уровня дикого типа.
	
	\subsection{Слой 3 – Протеостаз}
	Вероятность агрегации зависит от совместимости шаперон–клиент \(C_{AB}\). Оптимизация трёх шаперонов даёт общее снижение агрегации
	\[
	\frac{\mathrm{Agg}_{\mathrm{LL}}}{\mathrm{Agg}_{\mathrm{WT}}} \approx \sqrt[3]{0.78 \times 0.79 \times 0.78} \approx 0.78,
	\]
	т.е. примерно \(22\%\) уменьшение числа телец включения.
	
	\subsection{Слой 4 – Резонанс ВКМ–интегрин}
	Три оптимизированных компонента адгезии к внеклеточному матриксу дают
	\[
	\frac{\mathrm{Error}_{\mathrm{LL}}}{\mathrm{Error}_{\mathrm{WT}}} \approx \sqrt[3]{0.83 \times 0.83 \times 0.84} \approx 0.83.
	\]
	Скорость миграции обратно пропорциональна ошибкам адгезии. Поэтому мы ожидаем относительную скорость
	\[
	\frac{v_{\mathrm{LL}}}{v_{\mathrm{WT}}} \approx \frac{1}{0.83} \approx 1.20,
	\]
	т.е. увеличение на 20\%.
	
	\subsection{Комбинированное предсказание}
	Если предположить, что четыре слоя отказывают независимо, общая вероятность того, что клетка избежит фатальной ошибки, равна произведению вероятностей успеха каждого слоя. Следовательно, комбинированный фактор снижения фатальных ошибок (используя средние геометрические каждого слоя) составляет
	\[
	\frac{\varepsilon_{\mathrm{LL}}^{\mathrm{fatal}}}{\varepsilon_{\mathrm{WT}}^{\mathrm{fatal}}} \approx 0.86 \times 0.80 \times 0.78 \times 0.83 \approx 0.45.
	\]
	При гипотезе, что репликативная продолжительность жизни масштабируется как \(1/\varepsilon\), ожидаемое продление жизни равно
	\[
	\frac{\mathrm{Lifespan}_{\mathrm{LL}}}{\mathrm{Lifespan}_{\mathrm{WT}}} \approx \frac{1}{0.45} \approx 2.2.
	\]
	Это верхняя граница, пренебрегающая перекрёстными взаимодействиями; консервативная оценка, учитывающая перекрывающиеся пути повреждений, помещает фактор в диапазон \(1.5\)–\(2.0\).
	
	\section{Экспериментальный протокол}
	\label{sec:protocol}
	
	\subsection{Клеточная система}
	Кардиомиоциты, полученные из человеческих индуцированных плюрипотентных стволовых клеток (iPSC‑CMs, например Coriell GM25256). Клетки культивируют в RPMI 1640 + B27 при 37\(^\circ\)C, 5\% CO\(_2\).
	
	\subsection{Генетические конструкции}
	Для каждого гена из Таблицы~\ref{tab:targets} готовят две синтетические кДНК:
	\begin{itemize}
		\item \textbf{WT‑OE}: нативная кодирующая последовательность, клонированная в доксициклин-индуцируемый лентивирусный вектор (pLV‑TRE‑Tight‑IRES‑EGFP).
		\item \textbf{LL‑OPT}: оптимизированная BCLM последовательность в том же векторе.
	\end{itemize}
	Для BRCA1 создаётся скремблированный контроль (SCR) путём случайной перестановки синонимичных кодонов.
	
	\subsection{Контроли}
	\begin{enumerate}
		\item Пустой вектор (EV).
		\item WT‑OE с сопоставимыми уровнями белка.
		\item SCR (контроль специфичности).
	\end{enumerate}
	
	\subsection{График измерений}
	\begin{table}[H]
		\centering
		\caption{Экспериментальный график.}
		\label{tab:schedule}
		\begin{tabular}{@{}llll@{}}
			\toprule
			\textbf{День} & \textbf{Анализ} & \textbf{Слой} & \textbf{Показатель} \\
			\midrule
			0–7   & Трансдукция, селекция & --- & EGFP \\
			7–14  & Индукция доксициклином & Все & уровни белков \\
			14–21 & Анализ мутаций HPRT & 1 & колонии 6‑TG\(^{\mathrm{R}}\) \\
			14–21 & MitoSOX, TMRE & 2 & проточная цитометрия \\
			21–28 & Анализ включений HTT‑Q72 & 3 & флуоресцентная микроскопия \\
			21–28 & Скретч-анализ (рана) & 4 & скорость закрытия раны \\
			28–56 & Серийное пассирование & Все & сенесценс \(\beta\)-gal, теломерная qPCR \\
			\bottomrule
		\end{tabular}
	\end{table}
	
	\subsection{Репликативная продолжительность жизни}
	Клетки пассируют при 80\% конфлюэнтности; регистрируют кумулятивные удвоения популяции. Прогнозируется, что клетки BCLM‑LL‑OPT достигнут в \(1.5\)–\(2.0\) раза больше удвоений, чем контрольные клетки WT‑OE.
	
	\section{Контролируемая координационная регрессия (CCR)}
	\label{sec:CCR}
	Для доказательства причинности оптимизированные гены временно подавляются.
	
	\begin{table}[H]
		\centering
		\caption{CCR-вмешательства.}
		\label{tab:CCR}
		\begin{tabular}{@{}llc@{}}
			\toprule
			\textbf{Слой} & \textbf{Метод} & \textbf{Целевой \(\REL\) после CCR} \\
			\midrule
			1 & Dox-индуцируемая shRNA против BRCA1\_LL & 0.62 (уровень WT) \\
			2 & CRISPR-бейс-редактирование NDUFS1 (I182F)       & 0.55 \\
			3 & Tet‑CRISPRi HSPA1A\_LL               & 0.52 \\
			4 & siRNA ITGB3\_LL                      & 0.57 \\
			\bottomrule
		\end{tabular}
	\end{table}
	
	Через 72 часа биомаркеры должны сместиться в сторону WT фенотипа. После отмывки или разбавления siRNA клетки должны вернуться в состояние LL в течение 48–72 часов. Количественное предсказание:
	\[
	\frac{\varepsilon_{\mathrm{CCR}}}{\varepsilon_{\mathrm{LL}}} = \left(\frac{\REL_{\mathrm{LL}}}{\REL_{\mathrm{CCR}}}\right)^{2/3}.
	\]
	
	\section{Фальсифицируемость}
	\label{sec:falsifiability}
	Все предсказанные значения в Таблице~\ref{tab:predictions} предварительно зарегистрированы. Если после трёх независимых повторений экспериментальные результаты для четырёх из пяти анализов выйдут за пределы предсказанных диапазонов, координационная теория старения YUCT будет считаться опровергнутой.
	
	\begin{table}[H]
		\centering
		\caption{Фальсифицируемые предсказания.}
		\label{tab:predictions}
		\begin{tabular}{@{}llll@{}}
			\toprule
			\textbf{Предсказание} & \textbf{Анализ} & \textbf{Ожидаемое (LL/WT)} & \textbf{Опровержение, если} \\
			\midrule
			Частота мутаций       & HPRT           & \(0.86 \pm 0.05\) & LL/WT \(> 0.95\) \\
			ROS                 & MitoSOX        & \(0.80 \pm 0.05\) & LL/WT \(> 0.90\) \\
			Агрегаты          & HTT‑Q72 фокусы   & \(0.78 \pm 0.05\) & LL/WT \(> 0.90\) \\
			Скорость миграции     & Скретч-анализ  & \(1.20 \pm 0.05\) & LL/WT \(< 1.05\) \\
			Репликативная продолжительность жизни & Кумулятивные PD  & \(1.5\)–\(2.0\)    & LL/WT \(< 1.3\) \\
			\bottomrule
		\end{tabular}
	\end{table}
	
	\section{Этическое примечание}
	\label{sec:ethics}
	Данный протокол использует исключительно культивируемые клетки. Генетические конструкции не упаковываются в вирусные векторы, пригодные для \textit{in vivo} доставки. Этические и управленческие аспекты рассмотрены в Приложении~\(\Sigma\).
	
	\section{Заключение}
	Протокол BCLM‑LL предоставляет полную, самосогласованную экспериментальную проверку теории старения YUCT. Каждое численное предсказание выведено из универсального закона ошибок и вычисленных с помощью BCLM значений \(\REL\); никакие свободные параметры не подгоняются \emph{post hoc}. Все числа в тексте согласованы с лежащей в основе Таблицей~\ref{tab:targets}, что гарантирует полную внутреннюю непротиворечивость. Положительный результат станет первой демонстрацией того, что клеточное здоровое долголетие может быть продлено путём рационального восстановления эффективности координации до ювенильного уровня.
	
	\vspace{0.5cm}
	\noindent\textbf{Данные и код:} Все последовательности, вычисленные значения \(\REL\) и предварительно зарегистрированные предсказания доступны по адресу \url{https://github.com/Alexey-Yakushev-YUCT/YPSDC}.
	
	\begin{thebibliography}{99}
		\bibitem{BCLM} A.\,V.\,Yakushev, ``Приложение BCLM: Биологический лагранжиан координации,'' in \emph{YUCT}, Zenodo (2026).
		\bibitem{AppendixL} A.\,V.\,Yakushev, ``Приложение L: Фрактальное масштабирование ошибок координации,'' in \emph{YUCT}, Zenodo (2026).
		\bibitem{AppendixP} A.\,V.\,Yakushev, ``Приложение P: Температурная зависимость гравитации,'' in \emph{YUCT}, Zenodo (2026).
		\bibitem{AppendixSigma} A.\,V.\,Yakushev, ``Приложение \(\Sigma\): Этические императивы и управление технологиями на основе YUCT,'' in \emph{YUCT}, Zenodo (2026).
		\bibitem{Bjedov2021} I.~Bjedov \emph{et al.}, ``A single hyper‑accurate mutation in the ribosome extends lifespan in yeast, worms, and flies,'' \emph{Cell Metabolism}, 33, 1054–1070 (2021).
		\bibitem{Kerekes2025} K.~Kerekes \emph{et al.}, ``Codon optimisation and evolutionary pressure in long‑lived mammals,'' \emph{bioRxiv} (2025).
		\bibitem{YPSDC_repo} A.\,V.\,Yakushev, YPSDC репозиторий, \url{https://github.com/Alexey-Yakushev-YUCT/YPSDC}.
	\end{thebibliography}
	
\end{document}